% Generated by report/make_restated.py from theory/A3_lkatyusha_lyapunov.tex: do not edit (edit the note and rerun).
\begin{restated}{Remark}{rem:naive-book}{A3}{0}{relation with Nesterov's method of the book}
Eliminating $z$ from \eqref{eq:naive}: $y_{t+1}=\theta_1z_{t+1}+(1-\theta_1)y_t$ and $x_{t+1}=\theta_1z_{t+1}+(1-\theta_1)y_{t+1}$ give $x_{t+1}=y_{t+1}+(1-\theta_1)(y_{t+1}-y_t)$, which is the extrapolation (5.17) of the book, $\eta_t=\theta_t+\beta(\theta_t-\theta_{t-1})$ (a point, in the book's notation; it is written $v_t$ in the outcome of the experiments, Section~\ref{sec:lk-instrument}), with $\beta=1-\theta_1$. The book takes $\beta=\frac{1-\sqrt{\mu/L}}{1+\sqrt{\mu/L}}=1-2\sqrt{\mu/L}+\mathcal O(\mu/L)$: with $\theta_1=\sqrt{\mu/L}$ the two momentum coefficients agree at zeroth order (both are $1-\Theta(\sqrt{\mu/L})$), but the momentum deficits differ by a factor $2$ ($1-\beta=\sqrt{\mu/L}$ here, against $2\sqrt{\mu/L}+\mathcal O(\mu/L)$ in the book). The gradient step reads $y_{t+1}=x_t-\frac{\theta_1\eta}{L(1+\eta\sigma)}g_t+\frac{\theta_1\eta\sigma}{1+\eta\sigma}(x_t-z_t)$: it is (5.16) with $\Fp(\eta_{t-1})$ replaced by $g_t$, step $\frac{\theta_1\eta}{L(1+\eta\sigma)}$, and a drift which comes from putting the strong convexity in the mirror step. Since $x_t-z_t=-\frac{1-\theta_1}{\theta_1}(x_t-y_t)$, the drift equals $-\frac{\eta\sigma(1-\theta_1)}{1+\eta\sigma}(x_t-y_t)$, i.e.\ it is of relative size $\eta\sigma$ with respect to the momentum term $x_t-y_t=(1-\theta_1)(y_t-y_{t-1})$, where $\eta\sigma=\sqrt\sigma$ for the parameters of Proposition~\ref{prop:agd-lc} and $\eta\sigma=\frac{\sigma}{3\theta_1}$ for \eqref{eq:lk-param} ($=\sqrt{\sigma/(6n)}$ when $\kappa\ge8n/3$). The drift is therefore small compared with the momentum term, but of the same order $\mathcal O(\sqrt\sigma)$ as the discrepancy between the two momentum coefficients, which is the order that governs the accelerated rate: the two-sequence form of the book and \eqref{eq:naive} coincide only up to $\mathcal O(\sqrt\sigma)$ corrections of the momentum term. (Carrying the strong convexity by a proximal term does not remove the drift. If $f_i=\ell_i+\frac\mu2\norm{\cdot}^2$ with $\ell_i$ convex, as for $\ell_2$-regularized logistic regression with $\mu=\lambda$, then $g_t=\tilde g_t+\mu x_t$ with $\tilde g_t$ the SVRG estimator of $\frac1n\sum_i\ell_i$, and \eqref{eq:K3} is exactly $z_{t+1}=\mathrm{prox}_{\gamma g}(z_t-\gamma\tilde g_t)=\frac{z_t-\gamma\tilde g_t}{1+\gamma\mu}$ with $g=\frac\mu2\norm{\cdot}^2$ and $\gamma=\eta/L$, the $z$-step of Algorithm~1 of \cite{driggs2022accelerating}: the drift is the gradient of $g$ taken at $z_t$ instead of $x_t$. On such problems \eqref{eq:naive} is therefore their Algorithm~1 with $\tau=\theta_1$ and $\gamma=\eta/L$, except that their SVRG variant refreshes the reference point at the query point.) The two-sequence form of the book with $\Fp$ replaced by a mini-batch SVRG estimator, a proximal step, a reference point refreshed every $m$ iterations and the momentum restarted at each refresh ($x_1=y_1=\tilde x$ in their notation) is the method of Nitanda \cite{nitanda2014stochastic}; Allen-Zhu writes that it ``corresponds to this paper but \emph{without} Katyusha momentum (i.e., $\tau_2=0$)'' \cite[Appendix~A]{allenzhu2018katyusha}, and Driggs et al.\ write their scheme in the same two-sequence form (``$x_{k+1}=y_k+(1-\tau_k)(y_k-y_{k-1})$'') \cite[Section~2]{driggs2022accelerating}.
\end{restated}
